\insertImage{fde831f5-1dd5-468b-8682-1d0e38a05cc0.png}{Aquarelle d'une réunion d'équipe dans une entreprise libérée, avec des employés partageant des idées et prenant des décisions ensemble.}

\chapter{La communication et la prise de décision}

Faites l'exercice avec moi. Dans votre entreprise, aujourd'hui, qui décide d'embaucher quelqu'un ? Qui décide d'un achat à dix mille euros ? Qui décide qu'un projet s'arrête ?

Si vous avez répondu du tac au tac, vous travaillez dans une organisation classique : la réponse, c'est un nom ou un titre, et tout le monde le connaît. Dans une entreprise libérée, ces trois questions déclenchent une conversation. Ce n'est pas un bug, c'est le principe : les décisions descendent au plus près de ceux qui font, et l'information circule à l'horizontale plutôt que de remonter et redescendre les étages\footnote{Carney, B. M., \& Getz, I. (2009). Freedom, Inc.: Free Your Employees and Let Them Lead Your Business for Greater Success. Crown Business.}.

Sur le papier, c'est plus rapide, plus intelligent, plus engageant. Dans la vraie vie, c'est tout ça — et c'est aussi le premier endroit où une libération se met à grincer. Voyons pourquoi.

\section{La difficulté de trouver un équilibre entre autonomie et cohésion}

Chez FAVI, la fonderie picarde devenue l'icône française du modèle, les opérateurs géraient leur rythme de production, leurs délais, jusqu'à l'entretien de leurs machines\footnote{Zobrist, J.-F. (2012). La belle histoire de FAVI: L'entreprise qui croit que l'homme est bon. Humanisme \& Organisations.}. Résultat : de la réactivité, de la fierté, de l'innovation. L'autonomie tient ses promesses — ce n'est pas le sujet.

Le sujet, c'est ce qui se passe quand tout le monde est autonome en même temps.

Prenez Valve, le studio américain de jeux vidéo célèbre pour son organisation sans hiérarchie : chacun choisit les projets sur lesquels il travaille, littéralement en roulant son bureau jusqu'à l'équipe de son choix\footnote{Valve Corporation (2012). Handbook for New Employees.}. C'est séduisant. C'est aussi mathématique : si chacun va vers le projet le plus excitant, qui travaille sur le projet nécessaire mais ingrat ? L'autonomie individuelle ne produit pas spontanément une cohérence collective. Personne n'a tort, chacun optimise son coin — et l'ensemble se fragmente.

Même mécanique sur les objectifs : une équipe privilégie l'innovation, une autre la réduction des coûts, et les deux ont raison localement. Sans un alignement explicite, ces bonnes décisions locales s'additionnent en incohérence globale. C'est exactement le service invisible que rendait la hiérarchie — elle alignait, souvent mal, mais elle alignait. La supprimer sans la remplacer par autre chose, c'est garder le problème et perdre la solution.

Trop de cohésion étouffe l'initiative. Trop d'autonomie dissout le collectif. Toute entreprise libérée vit sur cette ligne de crête, en permanence.

\section{Les risques de désorganisation et de confusion}

Il faut raconter l'histoire de GitHub, parce qu'elle est arrivée aux gens les plus intelligents du monde — ce qui devrait tous nous rendre modestes.

Au début des années 2010, GitHub, plateforme adorée des développeurs, fonctionne sans managers : chacun choisit son rôle et ses projets. Puis l'entreprise grandit, et les questions sans réponse s'accumulent. Qui tranche quand deux équipes veulent des choses incompatibles ? À qui parle un nouveau qui ne comprend pas où va l'entreprise ? Qui s'occupe des carrières ? En 2014, après une crise interne retentissante, GitHub réintroduit une structure de management. Non pas parce que la hiérarchie serait morale ou sacrée — mais parce que certaines fonctions qu'elle assurait n'avaient jamais été réattribuées à personne\footnote{Foss, N. J., \& Klein, P. G. (2014). Why Managers Still Matter. MIT Sloan Management Review, 56(1), 73-80.}.

La leçon n'est pas « la libération ne marche pas ». La leçon, c'est qu'une organisation sans hiérarchie a \emph{plus} besoin de clarté qu'une organisation classique, pas moins : qui est responsable de quoi, comment on tranche un désaccord, où circule l'information. Le flou n'est pas de la liberté. Le flou, c'est du pouvoir informel — celui de ceux qui parlent le plus fort, et il est bien plus difficile à contester qu'un organigramme.

Autre risque, plus prosaïque : la lenteur. Quand chaque décision exige de consulter tout le monde, on ne décide plus, on réunionne. Automattic, l'entreprise entièrement distribuée derrière WordPress, a dû construire tout un outillage — blogs internes, coordination asynchrone — pour que des centaines de personnes réparties sur la planète restent alignées sans se noyer\footnote{Mullenweg, M. (2020). The Five Levels of Autonomy. Ma.tt.}. La coordination sans hiérarchie n'est pas gratuite. Elle se paie en process, en outils et en discipline — juste ailleurs.

\section{Les outils et méthodes pour faciliter la communication et la prise de décision}

La bonne nouvelle, c'est que ces problèmes ont des solutions éprouvées. La moins bonne, c'est qu'aucune ne s'installe en cochant une case.

Côté circulation de l'information, les entreprises libérées s'appuient massivement sur l'écrit et la transparence par défaut : des outils de gestion de projet partagés, et surtout une culture où l'information est publiée plutôt que transmise. Chez Automattic, les décisions se prennent et se documentent sur des blogs internes accessibles à tous — ce qui permet à n'importe qui de comprendre pourquoi un choix a été fait, six mois plus tard, sans avoir assisté à la réunion\footnote{Berkun, S. (2013). The Year Without Pants: WordPress.com and the Future of Work. Jossey-Bass.}.

Côté décision, arrêtons-nous sur une distinction que je trouve précieuse : le consensus et le consentement, que tout le monde confond. Le consensus demande que tout le monde soit \emph{pour} — c'est épuisant, et ça donne un droit de veto au plus têtu. Le consentement, popularisé par la sociocratie puis par l'holacratie, demande seulement que personne n'ait d'objection raisonnable et argumentée\footnote{Endenburg, G. (1998). Sociocracy: The organization of decision-making, "no objection!" as the principle of sociocracy. Eburon.}. La nuance a l'air subtile ; en pratique, c'est la différence entre des décisions qui aboutissent et des réunions qui recommencent. L'holacratie pousse la logique plus loin encore, avec des rôles explicites et des rituels de décision très codifiés\footnote{Robertson, B. J. (2015). Holacracy: The New Management System for a Rapidly Changing World. Henry Holt and Co.}.

Faut-il adopter l'un de ces systèmes clés en main ? Méfiez-vous de votre premier réflexe. Ces méthodes sont des points de départ, pas des destinations : chaque entreprise qui réussit finit par les tordre pour les adapter à sa culture et à son métier. Celle qui applique le manuel à la lettre a juste remplacé son ancienne bureaucratie par une bureaucratie plus exotique.

L'outil ne fait pas la liberté. Il la rend seulement possible — le reste, c'est vous.
